% ==============================================================================
% future_directions.tex
% ==============================================================================

\section{Future Directions}
\label{sec:future_directions}

Based on our analysis of the current literature, we identify several promising directions for future research:

\begin{enumerate}
  \item \textbf{Scalable Diffusion Models}: Developing efficient diffusion architectures (e.g., latent diffusion, consistency models) that can generate high-resolution 3D CT volumes in clinically acceptable time frames.

  \item \textbf{Closing the Domain Gap}: Creating domain adaptation and domain randomisation strategies that enable models trained on synthetic DRRs to generalise to real clinical X-rays without fine-tuning.

  \item \textbf{Standardised Benchmarks}: Establishing community-wide benchmarks with shared datasets, evaluation protocols, and standardised train/test splits to enable fair comparison across methods.

  \item \textbf{Hybrid Architectures}: Combining the strengths of multiple paradigms — for example, using neural fields for continuous representation with diffusion-based priors for improved generation quality.

  \item \textbf{Clinical Validation}: Moving beyond synthetic evaluation to prospective clinical studies that assess the diagnostic utility of reconstructed CT volumes.

  \item \textbf{Uncertainty Quantification}: Incorporating calibrated uncertainty estimates into reconstructions to flag regions where the model is least confident, which is critical for clinical decision-making.

  \item \textbf{Multi-Modal Integration}: Leveraging complementary information from other imaging modalities (e.g., ultrasound, MRI) alongside X-rays to improve reconstruction accuracy.
\end{enumerate}

% TODO: Expand each direction with supporting references and detailed discussion.
